
\section{Discussion: scope, and a reading of the results}\label{sec:discussion}

\subsection{Scope}\label{sec:scope}

The complexity results concern promise problems on explicitly presented
shared-probe models or budgeted signed-binomial systems in binary.
The decision bounds are for promised inputs. The algebraic endpoint test
checks feasibility of a supplied point. For algebraic Fine Connectivity,
boundedness of the rational feasible set remains a promise.
Theorems~\ref{thm:fixed-profile} and~\ref{thm:directed} rest on four
features of the model, whose roles Remark~\ref{rem:restrictions} gives: the
$\ell^1$ budget, equal word lengths, the probability floor
$\theta_m\ge\delta>0$, and register constraints that are affine and
separate from the rest.
The theorems are stated for the shared-probe model, its parity-mixture
kernels and its test families.

We claim no unconditional undecidability and do not classify test
languages: the boundary is a threshold at one syntactic step. The NCL construction uses calibrations, selector equations,
budgets and a growing probe family, with no necessity claimed for any of
them.
Restricting to rational points does not by itself create the obstructions.
The two signs after acquisition lie in different real components of the
binomial eigenvalue locus (Corollary~\ref{cor:real}). The algebraic NCL
reduction also gives obstructions over the reals
(Section~\ref{sec:real-ncl}). In Section~\ref{sec:real-rational}, however,
both endpoints are feasible rational configurations, and a real path joins
them even though rational fine reachability fails. The results concern a mathematical model and
make no claim about natural systems.

\subsection{A reading of the results}\label{sec:reading}

% lint: program begin
The question came from a program about how descriptions are strictly
extended, and the results admit a reading in its terms.
Foundations~I~\cite[Sections~4.1--4.2]{foundations1} takes the objects of a
theory to be the fixed points of a completion rule. Here, for a fixed
test list, the core is an interior operator on supports. Its
removed-coordinate map $Z\mapsto J\setminus\core(J\setminus Z)$ is a
closure operator satisfying Foundations~I's three order-closure axioms,
whose closed sets are the complements of the irreducible supports
(Proposition~\ref{prop:interior}). That much is proved; identifying these closed sets with the program's
objects is a reading.

Foundations~I calls the appearance of stable fixed
points within a theory package emergence and reserves open-endedness for
strict theory extension~\cite[Section~3.4]{foundations1}. We read feasible
revisions in which a configuration acquires a held test and retains its
earlier tests as an operational analogue of strict theory extension. An
acquisition can be redundant. The reading does not check that definability
changes strictly.
By Lemma~\ref{lem:restriction}, for single-word tests every terminal
peeling sequence from a given signed support has the same signed core.
As an analogy, not an instance, this matches the order-free admission of
Foundations~VII, Theorem~B~\cite{foundations7}.
The positive load equation in the calibrated NCL construction is analogous
to the disabling example in Foundations~VII. On the load line either
summand can vanish, and both cannot.
On explicitly presented models and
promised-feasible rational endpoints, single-word Fine Reachability is
P-complete under deterministic log-space many-one reductions and has
polynomial-size, polynomial-time checkable certificates for both answers.
With exactly one word-against-pair test, it is PSPACE-hard under
polynomial-time many-one reductions. Both statements hold at a fixed
profile and with profile steps allowed (Theorems~\ref{thm:boundary}
and~\ref{thm:directed}). For algebraic fine connectivity, with budgets and metric
fixed, adding tests refines the classes on points satisfying the enlarged
test list (Corollary~\ref{cor:refine}).
% gloss: adding tests refines the classes (budgets and metric fixed; points satisfying the enlarged list)

The reading has a one-line summary. A feasible world is what it cannot let
go of; all it keeps of that is a side; and the first commitment letting it
let go of either of two things, but not both, is the first wall. In the paper's terms, a
world is a feasible configuration and a commitment is a held test. The
signed core is what survives peeling, with its signs. In the calibrated
NCL construction, the fixed positive sum rules out joint vanishing, while
the selector equations and budgets govern the legal choices. Section~\ref{sec:acquisition} proves that early acquisition can
block fine reachability that late acquisition allows, even when both
endpoints pass the new test. We read this dependence on timing as a degree
of freedom in open-ended growth. None of this is a claim about nature, and
we make no priority claim within the program.
% lint: program end
